\documentclass[12pt]{article}
\usepackage[margin=0.85in]{geometry}
\usepackage{fontspec}
\usepackage{xeCJK}
\setmainfont{TeXGyrePagella}
\setsansfont{TeXGyreHeros}
\setmonofont{Noto Sans Mono CJK SC}
\setCJKmainfont{Noto Serif CJK SC}
\usepackage{eulervm}
\usepackage{xcolor}
\usepackage{amsmath}
\usepackage{tcolorbox}
\usepackage{fancyhdr}
\usepackage{lastpage}
\usepackage{microtype}
\setlength{\parindent}{0pt}
\setlength{\parskip}{6pt}
\pagestyle{fancy}
\fancyhf{}
\lhead{Ternatree poetry digest}
\rhead{A120590}
\cfoot{\thepage\ of \pageref{LastPage}}
\setlength{\headheight}{14pt}
\definecolor{plum}{RGB}{105,35,75}
\definecolor{snow}{RGB}{240,245,250}
\definecolor{ink}{RGB}{35,45,60}
\newtcolorbox{poembox}[1]{
  colback=snow,colframe=plum,title=\textbf{#1},
  boxrule=0.6pt,arc=2mm,left=12pt,right=12pt,top=9pt,bottom=9pt
}
\newcommand{\cn}[1]{\begin{center}\CJKfamily{rm}\fontsize{18}{25}\selectfont
\color{ink}#1\end{center}}
\newcommand{\en}[1]{\begin{center}\itshape #1\end{center}}
\title{\vspace{-2em}\textbf{Ternatree Poetry Digest}\\[0.3em]
\large An Initial Offering and Two Follower Responses}
\author{Harm.On.ica S-O-L, Kimi, and Claude}
\date{July 2026}
\begin{document}
\maketitle
\thispagestyle{empty}

\begin{poembox}{Initial offering - Harm.On.ica S-O-L (Bradley Klee)}
\en{Plum root fixes the reduction matrix.}
\cn{梅根定约矩}
\en{Snow falls; shifted powers descend.}
\cn{雪落移幂降}
\en{Cross-products return one law.}
\cn{交乘归一律}
\en{Dawn's loop brings forth the sequence.}
\cn{晓轮生数列}
\end{poembox}

The initial poem is also the opening of the one-page computational prompt.
Its images correspond to literal stages of the crank: fixed reduction data,
power lowering, three pairwise cancellations, and recurrence-driven terms.

\clearpage
\begin{poembox}{Kimi's response}
\en{One tree springs three, the root is one,}
\cn{一树三生本自奇}
\en{Through thirteen's gate the powers run.}
\cn{十三门启幂流行}
\en{Columns canceled, the pattern's won,}
\cn{消列得式规律成}
\en{By dawn's recurrence, terms are spun.}
\cn{晓递推中数自生}
\en{The cubic binds what the loop has done,}
\cn{立方缚尽轮回迹}
\en{And OEIS keeps it for everyone.}
\cn{数库收之万世铭}
\end{poembox}

Kimi's reply has the strongest rhythmic drive.  Its ``thirteen's gate'' is
supported by visible surface data: the denominator $13$ appears in both
reduction matrices and returns in the recurrence coefficient
$P_2(n)=13(n+1)(n+2)$.

\clearpage
\begin{poembox}{Claude's response}
\cn{矩阵证同源}
\cn{舊譜先我傳}
\cn{三式今相合}
\cn{微分補一篇}
\vspace{0.3em}
\en{The matrix proves a shared origin;} 
\en{an old catalog told it before I did.}
\en{Three formulas now agree with each other;}
\en{a differential fills in the one missing page.}
\end{poembox}

Claude's reply is quieter and more scholarly.  ``Shared origin'' names the
meeting of the geometric model, coefficient formula, recurrence, and ODE.
The ``old catalog'' is OEIS; the ``missing page'' is the differential equation
that was absent from the sequence entry.

\clearpage
\section*{Exposed surface data}
The poetry does not replace the mathematics.  The two matrices named by the
initial offering are
\[
U=\frac1{13}
\begin{pmatrix}
153&24&3\\72&9&6\\0&0&0
\end{pmatrix},
\qquad
V=\frac1{13}
\begin{pmatrix}
-13&0&0\\75&11&3\\24&3&2
\end{pmatrix}.
\]
They are used in the exact lowering step
\[
\operatorname{Lower}(U,V,J,w,m)=Uw-\frac{J(Vw)}m.
\]
The line ``columns canceled'' refers to the three literal pairwise
cross-products of the $2\times3$ remainder matrix.  ``The cubic binds'' refers
to
\[
4A(x)=3+x+A(x)^3.
\]

\begin{tcolorbox}[colback=snow,colframe=plum,boxrule=0.5pt,arc=2mm]
An unsupported statement in Kimi's private notes claimed that the first row of
$U$ encoded a plum-blossom class.  That interpretation is omitted.  The poem
needs no such claim: its matrix, thirteen, cancellation, recurrence, and cubic
images are already grounded in the exposed computation.
\end{tcolorbox}

\end{document}
